% ============================================
% MEM-11 (ZH)
% 中文翻译 — SPEC-935-R356
% ============================================
\section*{\textbf{MEM-11} --- 隔离}
\markboth{MEM-11 --- 隔离}{}

\noindent \textit{被沉默者的病区}

\noindent \textbf{1923—1945 —— 被沉默者的病区}

\vspace{0.5em}\noindent\rule{\textwidth}{0.5pt}\vspace{0.5em}

\noindent 第二次杀人之后，我被转到被沉默者的病区。

\noindent 那是走廊尽头。一个不出现在官方记录里的病区。待在那里的患者对医院来说不存在。他们是医院维持着活着的身体——勉强活着——以免不得不解释太多的死亡。

\noindent 病区有十二间牢房。每边六间。牢房很小——两平方米。混凝土床。没有床垫。没有毯子。地上一个洞用来排泄。

\noindent 患者进来，就再也不出去。

\noindent 我看见男人四十岁进来，四十五岁死去。蜷缩着。头发花白。牙齿烂掉，皮肤满是伤口。

\noindent 我们不说话。不是不想说——是忘记了怎么说。沉默变成了一种病。先是声音消失。然后是词语。然后是思想。

\noindent 怪物是唯一说话的。

\noindent 它填满了沉默。

\noindent \textit{“你不会变得像他们一样，”它说。}

\noindent \textit{“像什么样？”}

\noindent \textit{“空洞。”}

\noindent \textit{“我已经是空洞的。”}

\noindent \textit{“不。你有日记。你写。只要你在写，你就存在。他们——}\textit{它指着其他患者}\textit{——他们停止了书写。停止了说话。停止了存在。”}

\noindent \textit{“他们还活着。”}

\noindent \textit{“活着不是存在。活着是身体。存在是故事。如果没有人记得你，你就不存在。”}

\noindent \textit{“会有人记得吗？”}

\noindent \textit{“我会。”}

\noindent 它安静了片刻。

\noindent \textit{“还有读日记的人。”}

\vspace{0.5em}\noindent\rule{\textwidth}{0.5pt}\vspace{0.5em}

\noindent 冬夜里，寒冷从墙壁的缝隙钻进来，不走了。我蜷缩着睡，把日记抱在胸口，用身体暖它。

\noindent 那些夜里，怪物很安静。

\noindent 但我感觉到它在看。

\noindent 在等待。

\noindent 它知道我死在那里。它知道日记会和我一起被埋葬。而它知道，总有一天，会有人翻开。

\noindent 日记是它的蛋。

\noindent 我是那个巢。

\vspace{0.5em}
\noindent\textit{\textbf{↪ 路线:} \rota{MEM-10} \rotasep \rota{MEM-12}}
